% Unnummeriert, damit die sieben Teile des Musters die Nummern 1 bis 7 tragen.
\section*{Rahmen und Methode}
\addcontentsline{toc}{section}{Rahmen und Methode}
\markright{Rahmen und Methode}
\label{sec:methode}

\subsection*{Gegenstand und Frage}

Neste Oyj, Espoo, stellt erneuerbaren Diesel und nachhaltigen Flugkraftstoff (Segment
Renewable Products, RP) her und betreibt die Raffinerie Porvoo (Oil Products, OP) sowie ein
Tankstellennetz (Marketing \& Services). Der Bericht betrachtet die Gesch\"aftsjahre 2017
bis 2025 sowie das erste Halbjahr 2026 und entscheidet \"uber die Aktie zum Kurs von
\je{\Kurs}{EUR} (\KursDatum).

Die Frage wird ausdr\"ucklich \emph{nach Horizont getrennt} gestellt. Neste steht 2026 an
einer Stelle, an der drei Zeitskalen verschiedene Antworten geben k\"onnen: Die Marge des
ersten Halbjahrs 2026 ist die h\"ochste der eigenen Reihe (kurzer Horizont), der Ausbau in
Rotterdam soll 2027 in Betrieb gehen und die Investitionen danach senken (mittlerer
Horizont), und \"uber f\"unf bis f\"unfzehn Jahre entscheidet, was ein Eigent\"umer aus dem
Gesch\"aft tats\"achlich entnehmen kann (langer Horizont). Teil~\ref{sec:votum} gibt deshalb
ein Votum je Horizont.

\subsection*{Quellen}

Prim\"arquellen sind die Gesch\"aftsberichte 2019 bis 2025, der Financial Statements
Release 2018 (f\"ur das Jahr 2017) und der Halbjahresfinanzbericht 2026, jeweils in der
englischen Originalfassung und lokal abgelegt. Jede Seitenangabe bezeichnet die
\emph{gedruckte} Seite. In den Gesch\"aftsberichten stimmt sie mit der PDF-Seite
\"uberein, im Halbjahresbericht liegt sie eine Seite darunter, weil das Deckblatt
ungez\"ahlt ist. Sekund\"arquellen sind als solche gekennzeichnet: Kurse (Yahoo Finance),
Analystenkonsens und Vergleichsmultiples (stockanalysis.com), Meldungen der
F\"uhrungskr\"afte nach Art.~19 MAR (neste.com), Rating (neste.com) und die H\"urde (MSCI).

\subsection*{Wie die Zahlen in dieses Dokument kommen}

Die Berichte werden seitenweise zerlegt, die gedruckte Seitenzahl wird aus der Fu{\ss}zeile
gelesen und \"uber die Folge der Seiten gepr\"uft. Die Werte liest Programmcode aus dem
Textlayer der Kennzahlenseiten und Kapitalflussrechnungen, bei der Kapitalflussrechnung nach
der Spaltenstellung unter dem Jahreskopf, weil dort eine Notenspalte und eine Randspalte
Zahlen tragen, die keine Werte sind. Danach pr\"uft \texttt{scripts/pruefe\_seiten.py} jeden
gespeicherten Prim\"arwert gegen den Text der Seite, die er als Beleg nennt:
\GeprueftAnzahl{} von \GeprueftGesamt{} Werten sind so belegt; \ExternAnzahl{}
Sekund\"arwerte stehen getrennt und tragen keine Seite.

Diese Pr\"ufung erkennt eine falsch zugeordnete Zahl nicht, wenn die Zahl auf der Seite
steht. Dagegen steht die Kette \"uber die Berichte: Jede Kennzahlenseite druckt drei Jahre,
jede Kapitalflussrechnung zwei, und derselbe Wert muss in jedem Bericht gleich stehen. Die
Kette hat beim ersten Lauf drei Br\"uche gemeldet, alle aus einer Ursache: Die
Nachbarzelle \enquote{- of weighted average number of shares} war in den Verschuldungsgrad
und den ROACE geraten.
Die verbleibenden Abweichungen sind Neudarstellungen des Unternehmens und in
Abschnitt~\ref{sec:luecken} einzeln genannt.

\subsection*{Die drei Gr\"o{\ss}en, auf denen der Bericht ruht}

\textbf{Freier Zufluss.} Operativer Mittelzufluss, vermindert um Investitionen in Sach- und
immaterielle Anlagen und um die Tilgung von Leasingverbindlichkeiten. Neste weist als
\enquote{Free cash flow} den Mittelfluss vor Finanzierungst\"atigkeit aus
(\Mio{\FreierCfNeste}{EUR} f\"ur 2025\Q{GBXXV}{80}); der Bericht rechnet mit der eigenen
Gr\"o{\ss}e, weil nur sie den Leasingabfluss enth\"alt, den ein Eigent\"umer ebenso tragen
muss wie eine Investition.

\textbf{H\"urde.} \Pz{\Huerde} p.\,a.\ Das ist die Rendite des MSCI World (Gross Returns,
USD) \"uber zehn Jahre nach dem Factsheet zum 31.08.2026\QW{MSCI World Index, Index
Factsheet 31.08.2026}{quelle:msci}{19.09.2026}. Eine Anlage in Neste muss diese Rendite
schlagen, sonst ist der Indexfonds die bessere Wahl. Als Sensitivit\"at dient die Rendite
seit Auflage des Index (1987) von \Pz{\HuerdeLang}. Dieselben beiden H\"urden tragen die
Berichte zu GEA und den \enquote{Dual Champions}; die Schwellenkurse sind damit
berichts\"ubergreifend vergleichbar.

\textbf{Endwert.} Der Buchwert des Eigenkapitals zum 30.06.2026, \je{\EkJeAktieHj}{EUR} je
Aktie\Q{HJ}{3}, zusammen \Mio{\Buchwert}{EUR} auf \AktienMio{} Mio.\ Aktien.

\annahme{Alle Rechnungen sind Vorsteuerrechnungen auf Ebene des Unternehmens. Der Zufluss
ist der Zufluss an das Unternehmen, nicht der nach Quellen- und Abgeltungsteuer beim
Anleger verbleibende Betrag. Die H\"urde ist eine Bruttoindexrendite; der Vergleich bleibt
damit gleichartig.}

\annahme{Der freie Zufluss wird in jedem Szenario \emph{konstant} fortgeschrieben. Das ist
keine Prognose, sondern der Verzicht auf eine. Den fehlenden Wachstumsterm holt der Bericht
in Teil~\ref{sec:votum} als Umkehrung zur\"uck: Er rechnet aus, welches Wachstum der heutige
Kurs voraussetzt, und h\"alt es gegen die eigene Reihe.}

\annahme{Der W\"ahrungsunterschied zwischen H\"urde (USD) und Zufluss (EUR) bleibt
unber\"ucksichtigt. Eine W\"ahrungsbereinigung der H\"urde h\"atte eine eigene, nicht
belegte Annahme zur Kursentwicklung verlangt.}
